% ==================================================================
% id:       prob-001
% topic:    מבוא להסתברות
% subtopic: מרחב הסתברות אחיד
% solution: full
% verified: no
% ==================================================================
בסיום כנס מתמטי מכובד, 10 מתמטיקאיות קבעו להיפגש בסניף ארומה שבאוניברסיטת בן-גוריון. להפתעתן, התברר שיש באוניברסיטה שלושה סניפים כאלה, ולכן כל אחת מהן בחרה סניף באקראי והגיעה אליו.
\begin{enumerate}
\item מה ההסתברות שכל המתמטיקאיות הצליחו להיפגש כמתוכנן, כלומר הגיעו כולן לאותו סניף?
\item מה ההסתברות שלסניף בבניין 56 לא הגיעה אף אחת?
\item השתמשו בעקרון ההכלה וההפרדה, וחשבו את ההסתברות שלכל סניף תגיע לפחות מתמטיקאית אחת.
\end{enumerate}
\begin{solbox}
\unverified
מרחב המדגם הוא כל הדרכים שבהן 10 המתמטיקאיות בוחרות סניפים, $|\Omega|=3^{10}$, ומרחב ההסתברות אחיד.
\begin{enumerate}
\item נסמן $A$ = "כל המתמטיקאיות הגיעו לאותו סניף". $|A|=3$, כי כולן יכולות להגיע לסניף הראשון, לשני או לשלישי. לכן
\[
\PP(A)=\frac{|A|}{|\Omega|}=\frac3{3^{10}}=\ans{\frac1{3^9}}.
\]
\item נסמן $B$ = "לסניף בבניין 56 לא הגיעה אף אחת". אז $|B|=2^{10}$, כי לכל מתמטיקאית יש רק 2 אפשרויות ולא 3. לכן
\[
\PP(B)=\frac{2^{10}}{3^{10}}=\ans{\Big(\frac23\Big)^{10}}.
\]
\item נסמן $A_i$ = "לסניף ה-$i$ הגיעה לפחות מתמטיקאית אחת", $i=1,2,3$. אנחנו מחפשים את $\PP(A_1\cap A_2\cap A_3)$. לפי המשלים ודה־מורגן:
\[
\PP(A_1\cap A_2\cap A_3)=1-\PP(\comp{A_1}\cup\comp{A_2}\cup\comp{A_3}),
\]
ולפי עקרון ההכלה וההפרדה:
\[
\begin{split}
\PP(\comp{A_1}\cup\comp{A_2}\cup\comp{A_3})={}&\PP(\comp{A_1})+\PP(\comp{A_2})+\PP(\comp{A_3})\\
&-\PP(\comp{A_1}\cap\comp{A_2})-\PP(\comp{A_1}\cap\comp{A_3})-\PP(\comp{A_2}\cap\comp{A_3})\\
&+\PP(\comp{A_1}\cap\comp{A_2}\cap\comp{A_3}).
\end{split}
\]

\begin{itemize}
\item כל $\comp{A_i}$ פירושו שלסניף ה-$i$ לא הגיעה אף אחת. כמו בסעיף (ב), $\PP(\comp{A_i})=(\frac23)^{10}$.
\item כל $\comp{A_i}\cap\comp{A_j}$ פירושו שכולן הגיעו לסניף הנותר, ולכן $\PP(\comp{A_i}\cap\comp{A_j})=\frac1{3^{10}}$.
\item $\comp{A_1}\cap\comp{A_2}\cap\comp{A_3}=\emptyset$, ולכן הסתברותו $0$.
\end{itemize}
לכן
\[
\PP(A_1\cap A_2\cap A_3)=1-\Big(3\cdot\Big(\frac23\Big)^{10}-3\cdot\Big(\frac13\Big)^{10}\Big)=\ans{\frac{3^9-2^{10}+1}{3^9}}\approx0.948.
\]
\end{enumerate}
\end{solbox}
